\subsection{Sparse verification: exact rates and a decisive counterexample}\label{sec:security14:sparse}
The completed full-relation protocol verifies the
degree and source relations explicitly in its NP
circuit. A sparse-query replacement is not secured
by the algebraic fold dictionary alone. The
following quantitative statements identify precisely
which piece of a new proximity theorem is missing.

\begin{definition}[Enlarged masked Reed--Solomon codes]\label{def:security14:codes}
At fold level $j$ put
\[
 M_j=2^{m+s-j},\quad
 N_j=2^{m-j},\quad
 B_j=\lceil B/2^j\rceil,\quad
 D_j=N_j+B_j,
\]
and define the code
\[
 \mathcal C_j:=
 \operatorname{RS}(\Omega_j,D_j)
 =\{(P(x))_{x\in\Omega_j}:
          \deg P<D_j\}
 \subseteq\F^{M_j}.
\]
Its Hamming distance is the
ordinary number of unequal
evaluation coordinates, and
$\delta_j$ denotes its
relative minimum distance.
\end{definition}

\begin{lemma}[Exact dimension and minimum distance]\label{lem:security14:rs-distance}
For $1\le D_j\le M_j$,
the code $\mathcal C_j$ has
dimension $D_j$, minimum
Hamming distance
\[
 \boxed{d_j=M_j-D_j+1}
\]
and relative distance
$\delta_j=(M_j-D_j+1)/M_j$.
\end{lemma}
\begin{proof}
Evaluation at $M_j$
distinct field points
is injective on the
space of polynomials of
degree below $D_j$,
because a nonzero polynomial
of degree at most
$D_j-1<M_j$ cannot vanish
at all domain points.
Thus the dimension is
$D_j$. The difference
of two distinct codewords
is the evaluation of a
nonzero polynomial of degree
at most $D_j-1$, with at most
$D_j-1$ roots on $\Omega_j$.
Its Hamming weight is
therefore at least
$M_j-(D_j-1)$.
Equality is attained by
a polynomial which is the
product of $D_j-1$ distinct
linear factors at evaluation
points (use the constant
polynomial $1$ when
$D_j=1$). Hence the
minimum distance is
exactly $M_j-D_j+1$.
\end{proof}

\begin{lemma}[Degree-rate bounds and required blow-up]\label{lem:security14:rates}
For $1\le B\le N$ and
$0\le j\le m$,
\[
 \boxed{
 \frac{D_j}{M_j}
 =2^{-s}\left(
 1+\frac{\lceil B/2^j\rceil}{N_j}
 \right)\le2^{1-s}.}
\]
In particular, $s\ge2$
guarantees rate at most
$1/2$ at every level.
When $s=1,B=N$, the
code is the entire ambient
space at every level.
\end{lemma}
\begin{proof}
Divide $D_j=N_j+B_j$
by $M_j=2^sN_j$.
Since $B\le N=2^m$,
we have
$B_j=\lceil B/2^j\rceil
\le\lceil2^{m-j}\rceil=N_j$.
This gives the inequality.
For $s=1,B=N$,
$B_j=N_j$ and hence
$D_j=2N_j=M_j$,
so evaluation polynomials
of degree below $M_j$
give every word in
$\F^{M_j}$ by
interpolation.
\end{proof}

\begin{lemma}[Folding a monomial halves its exponent]\label{lem:security14:half-exp}
For every integer $e\ge0$
and fold parameter $r\in\F$,
the polynomial $\pfold_r(X^e)$
has degree at most
$\lfloor e/2\rfloor$.
Consequently for any
$m$ fold parameters,
\[
 \boxed{
 \deg(\pfold_{r_{m-1}}\circ\cdots
 \circ\pfold_{r_0}(X^e))
 \le \lfloor e/2^m\rfloor.}
\]
\end{lemma}
\begin{proof}
The defining monomial
rules of
Lemmas~\ref{lem:kbff:fold-even}
and~\ref{lem:kbff:fold-odd}
send $X^{2a}$ and
$X^{2a+1}$ to
public scalar multiples
of $X^a$, respectively.
The scalar may vanish,
which can only reduce
the output degree.
Thus the degree after
one fold is at most
$\lfloor e/2\rfloor$.
Iterate this inequality
$m$ times and use
$\lfloor\lfloor e/2\rfloor/2
\rfloor=\lfloor e/4\rfloor$.
\end{proof}

\begin{theorem}[Sparse local folds and final linearity do not certify initial degree]\label{thm:security14:sparse-counterexample}
Assume $1\le B<N$ and
$M\ge2N$.
Let
\[
 \boxed{P_0(X):=X^{N+B}.}
\]
The evaluation word of $P_0$
is \emph{not} in
$\mathcal C_0=\operatorname{RS}
(\Omega_0,N+B)$.
Nevertheless, for \emph{every}
sequence of fold challenges
$r_0,\ldots,r_{m-1}\in\F$,
the chain
$P_{j+1}=\pfold_{r_j}(P_j)$
passes \emph{every}
polynomial/local-word fold
consistency check, and its
final polynomial satisfies
$\deg P_m<2$.
Therefore even exhaustive
local fold consistency
and final degree-one closure
do not imply the required
initial Reed--Solomon
membership.
\end{theorem}
\begin{proof}
The polynomial $P_0$ has
degree $N+B=D_0$.
Because $B<N$,
\[
 D_0=N+B<2N\le M.
\]
If its evaluation word
were in $\mathcal C_0$,
some polynomial $Q$
of degree below $D_0$
would agree with $P_0$
at all $M$ domain points.
Then $P_0-Q$ would be
nonzero of degree $D_0<M$
with $M$ distinct roots,
contradicting the univariate
root bound. Thus its word
is outside $\mathcal C_0$.

Define all later polynomials
by the \emph{actual}
normalized folding rule.
For each pair
$x,-x$ and square
$y=x^2$, the evaluation
fold identity
Theorem~\ref{thm:kbff:wordfold}
holds identically.
Hence every local
consistency check passes,
regardless of the chosen
challenges or the number
of queries.
By
Lemma~\ref{lem:security14:half-exp},
\[
 \deg P_m
 \le\left\lfloor
      \frac{N+B}{N}
    \right\rfloor=1,
\]
since $0<B<N$.
Thus the final degree-one
check also passes.
This proves the stated
non-implication.
\end{proof}

\begin{remark}[Precise interpretation of the counterexample]
The word $P_0$ above does
\emph{not} satisfy the
full protocol's preprocessing
source/mask relation.
The theorem is a
counterexample to the
\emph{proposed sparse
replacement of that relation
by local folding and a
final line check}, not a
forgery against the fully
checked Chapter~13 protocol.
An optimized sparse protocol
must enforce both
initial source/degree binding
and robust proximity of
the adversarial word chain.
\end{remark}

\begin{lemma}[Exact fixed-discrepancy sampling error]\label{lem:security14:query-detection}
Fix two words on a domain of
size $M_j$ which disagree
at exactly $e$ coordinates,
and sample $q$ positions
uniformly and independently
\emph{with replacement}
after both words are fixed.
The probability that all
sampled coordinates miss
the discrepancy is
\[
 \boxed{(1-e/M_j)^q.}
\]
\end{lemma}
\begin{proof}
The complement of the
discrepancy set has
$M_j-e$ points.
Each independent uniform
sample avoids the
discrepancy with probability
$(M_j-e)/M_j$.
The probability of all
$q$ misses is the
product of those
$q$ equal factors.
\end{proof}

\begin{definition}[Actual sparse proximity contract]\label{def:security14:proximity-contract}
A valid sparse KBFold security
theorem must prove a quantitative
bound $\varepsilon_{\rm prox}$
for the following full
adaptive experiment:
all committed word shells and
their associated codeword
commitments are fixed before
their corresponding query
indices; the adversary may
supply arbitrary non-codeword
words, may correlate words
across rounds, and may adapt
them to preceding public fold
challenges, but not to the
fresh query coins after a
root is bound.
The verifier's local fold
checks, initial source/degree
test, final line test and
claim consistency must be
specified. The reduction
must bound acceptance of
a chain not representing
one valid initial
degree-$<D_0$ masked word
and its legal folds.
The bound must use the
\emph{actual} budgets
$D_j/M_j$ from
Lemma~\ref{lem:security14:rates}.
\end{definition}

\begin{proposition}[A fixed discrepancy bound is not a proximity theorem]\label{prop:security14:not-proximity}
Lemma~\ref{lem:security14:query-detection}
does not imply
Definition~\ref{def:security14:proximity-contract}
without an additional
lower bound on discrepancy
for every invalid adversarial
chain.
\end{proposition}
\begin{proof}
The fixed-discrepancy
lemma assumes that the
two words and the number
$e$ of their disagreements
are already given.
A proximity theorem must
derive a positive
discrepancy lower bound
from non-membership
in the appropriate
degree-bounded code
under correlated adaptive
folds. The polynomial
in
Theorem~\ref{thm:security14:sparse-counterexample}
is outside the initial
code but has \emph{zero}
local fold discrepancy
at every location.
Thus local disagreement
cannot be inferred from
initial non-membership
alone. The missing
implication requires
separate code/commitment
structure.
\end{proof}

\begin{theorem}[Conditional sparse soundness composition]\label{thm:security14:sparse-sound}
Suppose a particular
sparse replacement of
Chapter~13's full
word-relation checks
satisfies
Definition~\ref{def:security14:proximity-contract}
with error
$\varepsilon_{\rm prox}$,
and retains the
prefix-preserving source
extraction, Pedersen binding,
Merkle collision resistance
and remaining valid
relation proofs of
Theorem~\ref{thm:security14:full-binding}.
Then its false-evaluation
acceptance probability is
at most
\[
 \boxed{\min\{1,\,
 \varepsilon_{\rm KE}
 +\varepsilon_{\rm DL}
 +\varepsilon_{\rm CRH}
 +\varepsilon_{\rm prox}\}.}
\]
\end{theorem}
\begin{proof}
As in
Theorem~\ref{thm:security14:full-binding},
exclude preprocessing
extraction failures,
conflicting registered
field openings and
conflicting authenticated
Merkle shells.
Their union probability
is bounded by the first
three errors. Outside
these events, the sparse
proximity contract ensures
except with probability
$\varepsilon_{\rm prox}$
that all accepted sparse
checks correspond to a
valid initial source-bound
masked word and its legal
fold chain. On that good
event the initial claimed
evaluation is the true
multilinear evaluation
of the extracted source
table by the same
exact algebraic argument
as in
Lemma~\ref{lem:full:fixed-eval}.
Thus false acceptance
belongs to the union
of the four bad events.
Apply the union bound.
\end{proof}

\begin{remark}[Locked research boundary]
No value of
$\varepsilon_{\rm prox}$
is supplied or implied
by the current Chapter~13
reference implementation.
The sparse optimization
is not declared secure.
Chapter~15 records this
unresolved component
inside its final
certification ledger;
no extra research
chapter is added.
\end{remark}
